\documentclass[10pt,a4paper]{amsart}
\usepackage{amsmath,amssymb,amsthm,graphicx}
\usepackage{verbatim}
\newenvironment{reference}{\comment}{\endcomment}
%\newenvironment{reference}{}{}
\newenvironment{slogan}{\comment}{\endcomment}
\newenvironment{history}{\comment}{\endcomment}

% For commutative diagrams we use Xy-pic
\usepackage[all]{xy}

% We use 2cell for 2-commutative diagrams.
\xyoption{2cell}
% Theorem environments.
%
\theoremstyle{plain}
\newtheorem{theorem}[subsection]{Theorem}
\newtheorem{proposition}[subsection]{Proposition}
\newtheorem{lemma}[subsection]{Lemma}

\theoremstyle{definition}
\newtheorem{definition}[subsection]{Definition}
\newtheorem{example}[subsection]{Example}
\newtheorem{exercise}[subsection]{Exercise}
\newtheorem{situation}[subsection]{Situation}

\theoremstyle{remark}
\newtheorem{remark}[subsection]{Remark}
\newtheorem{remarks}[subsection]{Remarks}

\numberwithin{equation}{subsection}

% Macros
%
\def\lim{\mathop{\mathrm{lim}}\nolimits}
\def\colim{\mathop{\mathrm{colim}}\nolimits}
\def\Spec{\mathop{\mathrm{Spec}}}
\def\Hom{\mathop{\mathrm{Hom}}\nolimits}
\def\Ext{\mathop{\mathrm{Ext}}\nolimits}
\def\SheafHom{\mathop{\mathcal{H}\!\mathit{om}}\nolimits}
\def\SheafExt{\mathop{\mathcal{E}\!\mathit{xt}}\nolimits}
\def\Sch{\mathit{Sch}}
\def\Mor{\mathop{\mathrm{Mor}}\nolimits}
\def\Ob{\mathop{\mathrm{Ob}}\nolimits}
\def\Sh{\mathop{\mathit{Sh}}\nolimits}
\def\NL{\mathop{N\!L}\nolimits}
\def\CH{\mathop{\mathrm{CH}}\nolimits}
\def\proetale{{pro\text{-}\acute{e}tale}}
\def\etale{{\acute{e}tale}}
\def\QCoh{\mathit{QCoh}}
\def\Ker{\mathop{\mathrm{Ker}}}
\def\Im{\mathop{\mathrm{Im}}}
\def\Coker{\mathop{\mathrm{Coker}}}
\def\Coim{\mathop{\mathrm{Coim}}}

% Boxtimes
%
\DeclareMathSymbol{\boxtimes}{\mathbin}{AMSa}{"02}

%
% Macros for moduli stacks/spaces
%
\def\QCohstack{\mathcal{QC}\!\mathit{oh}}
\def\Cohstack{\mathcal{C}\!\mathit{oh}}
\def\Spacesstack{\mathcal{S}\!\mathit{paces}}
\def\Quotfunctor{\mathrm{Quot}}
\def\Hilbfunctor{\mathrm{Hilb}}
\def\Curvesstack{\mathcal{C}\!\mathit{urves}}
\def\Polarizedstack{\mathcal{P}\!\mathit{olarized}}
\def\Complexesstack{\mathcal{C}\!\mathit{omplexes}}
% \Pic is the operator that assigns to X its picard group, usage \Pic(X)
% \Picardstack_{X/B} denotes the Picard stack of X over B
% \Picardfunctor_{X/B} denotes the Picard functor of X over B
\def\Pic{\mathop{\mathrm{Pic}}\nolimits}
\def\Picardstack{\mathcal{P}\!\mathit{ic}}
\def\Picardfunctor{\mathrm{Pic}}
\def\Deformationcategory{\mathcal{D}\!\mathit{ef}}
\usepackage{cleveref}
\crefname{originaldefinition}{original definition}{original definitions}
\Crefname{originaldefinition}{Original definition}{Original definitions}
\makeatletter
\newlabel{algebra-definition-complete}{{0317}{source}{}{sourcepointer-1}{}}
\newlabel{algebra-definition-complete@cref}{{[originaldefinition][0][]0317}{[1][0][]source}{}{}{}}
\makeatother
\crefname{originallemma}{original lemma}{original lemmas}
\Crefname{originallemma}{Original lemma}{Original lemmas}
\makeatletter
\newlabel{algebra-lemma-hathat-finitely-generated}{{05GG}{source}{}{sourcepointer-2}{}}
\newlabel{algebra-lemma-hathat-finitely-generated@cref}{{[originallemma][0][]05GG}{[1][0][]source}{}{}{}}
\makeatother
\crefname{originallemma}{original lemma}{original lemmas}
\Crefname{originallemma}{Original lemma}{Original lemmas}
\makeatletter
\newlabel{algebra-lemma-limit-complete}{{09B8}{source}{}{sourcepointer-3}{}}
\newlabel{algebra-lemma-limit-complete@cref}{{[originallemma][0][]09B8}{[1][0][]source}{}{}{}}
\makeatother
\crefname{originaldefinition}{original definition}{original definitions}
\Crefname{originaldefinition}{Original definition}{Original definitions}
\makeatletter
\newlabel{formal-spaces-definition-affine-formal-spectrum}{{0AIF}{source}{}{sourcepointer-4}{}}
\newlabel{formal-spaces-definition-affine-formal-spectrum@cref}{{[originaldefinition][0][]0AIF}{[1][0][]source}{}{}{}}
\makeatother
\crefname{originalreference}{original reference}{original references}
\Crefname{originalreference}{Original reference}{Original references}
\makeatletter
\newlabel{lemma-closed-is-complete}{{0HB3}{source}{}{sourcepointer-5}{}}
\newlabel{lemma-closed-is-complete@cref}{{[originalreference][0][]0HB3}{[1][0][]source}{}{}{}}
\makeatother
\crefname{originalsection}{original section}{original sections}
\Crefname{originalsection}{Original section}{Original sections}
\makeatletter
\newlabel{section-non-abelian}{{07JQ}{source}{}{sourcepointer-6}{}}
\newlabel{section-non-abelian@cref}{{[originalsection][0][]07JQ}{[1][0][]source}{}{}{}}
\makeatother
\crefname{originallemma}{original lemma}{original lemmas}
\Crefname{originallemma}{Original lemma}{Original lemmas}
\makeatletter
\newlabel{stacks-lemma-when-stack-in-sets}{{0430}{source}{}{sourcepointer-7}{}}
\newlabel{stacks-lemma-when-stack-in-sets@cref}{{[originallemma][0][]0430}{[1][0][]source}{}{}{}}
\makeatother
\crefname{originallemma}{original lemma}{original lemmas}
\Crefname{originallemma}{Original lemma}{Original lemmas}
\makeatletter
\newlabel{stacks-sheaves-lemma-quasi-coherent}{{06WK}{source}{}{sourcepointer-8}{}}
\newlabel{stacks-sheaves-lemma-quasi-coherent@cref}{{[originallemma][0][]06WK}{[1][0][]source}{}{}{}}
\makeatother
\crefname{originalsection}{original section}{original sections}
\Crefname{originalsection}{Original section}{Original sections}
\makeatletter
\newlabel{stacks-sheaves-section-quasi-coherent}{{06WF}{source}{}{sourcepointer-9}{}}
\newlabel{stacks-sheaves-section-quasi-coherent@cref}{{[originalsection][0][]06WF}{[1][0][]source}{}{}{}}
\makeatother
\crefname{originalsection}{original section}{original sections}
\Crefname{originalsection}{Original section}{Original sections}
\makeatletter
\newlabel{stacks-sheaves-section-quasi-coherent-algebraic-stacks}{{06WU}{source}{}{sourcepointer-10}{}}
\newlabel{stacks-sheaves-section-quasi-coherent-algebraic-stacks@cref}{{[originalsection][0][]06WU}{[1][0][]source}{}{}{}}
\makeatother
\title{Quasi-coherent modules on a Noetherian affine formal algebraic space need not form an abelian category}
\author{The Stacks Project Authors}
\date{stacks-project:examples.tex:local-snapshot-20261012; GNU FDL 1.2 or later}
\begin{document}
\maketitle
\noindent\textit{Editorial note.} All delivered mathematical text is the author's own, copied from the pinned Stacks Project chapter file examples.tex; no statement or proof step was rewritten, reordered or removed.\par\medskip
\noindent\textit{Editorial note.} The counted proof is the author's complete argument: The author's footnote on the definition of quasi-coherent is part of the counted statement and is kept verbatim. The final 'this category is not abelian' of the context points at Section non-abelian, whose kernel/cokernel machinery is included as the first context span; the witnessing non-exactness is the delivered card L48306.\par\medskip
\noindent\textit{Editorial note.} The label lemma-quasi-coherent-not-abelian names the counted statement, and the proof environment immediately following it in the source is delivered verbatim as its proof.\par\medskip
\noindent\textit{Editorial note.} The counted statement is carried by the author's own lemma environment in the chapter source.\par\medskip
\noindent\textit{Allowed standard prerequisites (editorial).}\begin{itemize}
\item Context span 792-805 (section-non-abelian): for R a ring and I finitely generated, the category A of I-adically complete R-modules has cokernel of phi: M -> N given by the I-adic completion of the usual cokernel (Algebra, hathat-finitely-generated) and kernel the usual kernel, which is closed and I-adically complete by Lemma closed-is-complete (delivered card L48284)
\item Context span 1161-1201 (section-nonabelian-QCoh): X = Spf(Z\_p) with the p-adic topology is a sheaf in sets on (Sch/Z)\_fppf and gives a stack in setoids (Stacks, when-stack-in-sets); quasi-coherent modules on it correspond to systems (M\_n) of abelian groups annihilated by p\textasciicircum{}n with M\_n = M\_\{n+1\}/p\textasciicircum{}n M\_\{n+1\}, equivalently to the p-adically complete abelian group M = lim M\_n with M\_n = M/p\textasciicircum{}n M (Algebra, limit-complete); the category of quasi-coherent modules on X is thus equivalent to the category of p-adically complete abelian groups, which the context states is not abelian, see Section non-abelian
\end{itemize}\medskip
\ifdefined\hypertarget\hypertarget{sourcepointer-1}{}\fi
\noindent\textit{Source reference.} Original definition 0317: sources/stacks-chow/tags.txt, source lines 3974--3974. Prerequisite pointer invoked at 1222; it names a standard fact of a sibling chapter that the counted argument uses, and every step of the counted proof built on it is retained verbatim in the delivered body.\par\medskip
\ifdefined\hypertarget\hypertarget{sourcepointer-2}{}\fi
\noindent\textit{Source reference.} Original lemma 05GG: sources/stacks-chow/tags.txt, source lines 7060--7060. Prerequisite pointer invoked at 1227; it names a standard fact of a sibling chapter that the counted argument uses, and every step of the counted proof built on it is retained verbatim in the delivered body.\par\medskip
\ifdefined\hypertarget\hypertarget{sourcepointer-3}{}\fi
\noindent\textit{Source reference.} Original lemma 09B8: sources/stacks-chow/tags.txt, source lines 11786--11786. Prerequisite pointer invoked at 1268; it names a standard fact of a sibling chapter that the counted argument uses, and every step of the counted proof built on it is retained verbatim in the delivered body.\par\medskip
\ifdefined\hypertarget\hypertarget{sourcepointer-4}{}\fi
\noindent\textit{Source reference.} Original definition 0AIF: sources/stacks-chow/tags.txt, source lines 13263--13263. Prerequisite pointer invoked at 1244; it names a standard fact of a sibling chapter that the counted argument uses, and every step of the counted proof built on it is retained verbatim in the delivered body.\par\medskip
\ifdefined\hypertarget\hypertarget{sourcepointer-5}{}\fi
\noindent\textit{Source reference.} Original reference 0HB3: sources/stacks-examples/pinned-original-source.tex, source lines 808--808. Prerequisite pointer invoked at 1229; it names a standard fact of the same chapter that the counted argument uses, and every step of the counted proof built on it is retained verbatim in the delivered body.\par\medskip
\ifdefined\hypertarget\hypertarget{sourcepointer-6}{}\fi
\noindent\textit{Source reference.} Original section 07JQ: sources/stacks-examples/pinned-original-source.tex, source lines 790--790. Prerequisite pointer invoked at 1272; it names a standard fact of the same chapter that the counted argument uses, and every step of the counted proof built on it is retained verbatim in the delivered body.\par\medskip
\ifdefined\hypertarget\hypertarget{sourcepointer-7}{}\fi
\noindent\textit{Source reference.} Original lemma 0430: sources/stacks-chow/tags.txt, source lines 5347--5347. Prerequisite pointer invoked at 1247; it names a standard fact of a sibling chapter that the counted argument uses, and every step of the counted proof built on it is retained verbatim in the delivered body.\par\medskip
\ifdefined\hypertarget\hypertarget{sourcepointer-8}{}\fi
\noindent\textit{Source reference.} Original lemma 06WK: sources/stacks-chow/tags.txt, source lines 8819--8819. Prerequisite pointer invoked at 1255; it names a standard fact of a sibling chapter that the counted argument uses, and every step of the counted proof built on it is retained verbatim in the delivered body.\par\medskip
\ifdefined\hypertarget\hypertarget{sourcepointer-9}{}\fi
\noindent\textit{Source reference.} Original section 06WF: sources/stacks-chow/tags.txt, source lines 8814--8814. Prerequisite pointer invoked at 1232; it names a standard fact of a sibling chapter that the counted argument uses, and every step of the counted proof built on it is retained verbatim in the delivered body.\par\medskip
\ifdefined\hypertarget\hypertarget{sourcepointer-10}{}\fi
\noindent\textit{Source reference.} Original section 06WU: sources/stacks-chow/tags.txt, source lines 8828--8828. Prerequisite pointer invoked at 1236; it names a standard fact of a sibling chapter that the counted argument uses, and every step of the counted proof built on it is retained verbatim in the delivered body.\par\medskip
\noindent
Let $R$ be a ring and let $I \subset R$ be a finitely generated ideal.
Consider the category $\mathcal{A}$ of $I$-adically complete
$R$-modules, see
Algebra, Definition \ref{algebra-definition-complete}.
Let $\varphi : M \to N$ be a morphism of $\mathcal{A}$.
The cokernel of $\varphi$ in $\mathcal{A}$ is the $I$-adic completion
$(\Coker(\varphi))^\wedge$ of the usual cokernel
(as $I$ is finitely generated this completion is complete, see
Algebra, Lemma \ref{algebra-lemma-hathat-finitely-generated}).
Let $K = \Ker(\varphi)$; since $\varphi$ is continuous this
is a closed submodule of $M$. By Lemma \ref{lemma-closed-is-complete}
below we see that $K$ is $I$-adically complete and hence is the
kernel of $\varphi$ in $\mathcal{A}$.

In Sheaves on Stacks, Section \ref{stacks-sheaves-section-quasi-coherent}
we defined the category of quasi-coherent modules on a category fibred in
groupoids over $\Sch$. Although we show in
Sheaves on Stacks, Section
\ref{stacks-sheaves-section-quasi-coherent-algebraic-stacks}
that this category is abelian for algebraic stacks, in this
section we show that this is not the case for formal algebraic spaces.

\medskip\noindent
Namely, consider $\mathbf{Z}_p$ viewed as topological ring using
the $p$-adic topology. Let $X = \text{Spf}(\mathbf{Z}_p)$, see
Formal Spaces, Definition
\ref{formal-spaces-definition-affine-formal-spectrum}.
Then $X$ is a sheaf in sets on $(\Sch/\mathbf{Z})_{fppf}$
and gives rise to a stack in setoids $\mathcal{X}$, see 
Stacks, Lemma \ref{stacks-lemma-when-stack-in-sets}.
Thus the discussion of Sheaves on Stacks, Section
\ref{stacks-sheaves-section-quasi-coherent-algebraic-stacks}
applies. 

\medskip\noindent
Let $\mathcal{F}$ be a quasi-coherent module on $\mathcal{X}$.
Since $X = \colim \Spec(\mathbf{Z}/p^n\mathbf{Z})$ it is clear from
Sheaves on Stacks, Lemma \ref{stacks-sheaves-lemma-quasi-coherent}
that $\mathcal{F}$ is given by a sequence $(\mathcal{F}_n)$ where
\begin{enumerate}
\item $\mathcal{F}_n$ is a quasi-coherent module on
$\Spec(\mathbf{Z}/p^n\mathbf{Z})$, and
\item the transition maps give isomorphisms
$\mathcal{F}_n = \mathcal{F}_{n + 1}/p^n\mathcal{F}_{n + 1}$.
\end{enumerate}
Converting into modules we see that $\mathcal{F}$ corresponds to a
system $(M_n)$ where each $M_n$ is an abelian group annihilated
by $p^n$ and the transition maps induce isomorphisms
$M_n = M_{n + 1}/p^n M_{n + 1}$. In this situation the module
$M = \lim M_n$ is a $p$-adically complete module and $M_n = M/p^n M$, see
Algebra, Lemma \ref{algebra-lemma-limit-complete}.
We conclude that the category of quasi-coherent modules on $X$
is equivalent to the category of $p$-adically complete
abelian groups. This category is not abelian, see
Section \ref{section-non-abelian}.

\begin{lemma}
\label{lemma-quasi-coherent-not-abelian}
The category of quasi-coherent\footnote{With quasi-coherent modules
as defined above. Due to how things are setup in the Stacks project,
this is really the correct definition; as seen above our definition
agrees with what one would naively have defined to be quasi-coherent modules
on $\text{Spf}(A)$, namely complete $A$-modules.}
modules on a formal algebraic space
$X$ is not abelian in general, even if $X$ is a Noetherian affine
formal algebraic space.
\end{lemma}

\begin{proof}
See discussion above.
\end{proof}

\end{document}
